% Part: proof-theory
% Chapter: normalization
% Section: permutations

\documentclass[../../../include/open-logic-section]{subfiles}

\begin{document}

\olfileid[hi]{pt}{nor}{per}

\olsection{क्रम-विनिमय रूपांतरण}

\Log{N1i} में कट ऐसे खंड हैं जो किसी !!a{elimination} नियम के मुख्य
आधार-वाक्य पर समाप्त होते हैं। प्रसामान्यीकरण प्रमाण की एक स्थिति में महत्तम
कटों की लंबाई, और उसके साथ !!{proof} की कट लंबाई, घटानी होती है। लंबाई $>1$
वाला कट अनेक ढंगों से समाप्त हो सकता है, जो इस पर निर्भर है कि कट में~$!A$ की
अंतिम !!{formula}-आवृत्ति (\Elim{\lor} या \Elim{\lexists}) का निष्कर्ष है या
नहीं, और वह किस !!{elimination} अनुमान का मुख्य आधार-वाक्य है।

उदाहरणतः यदि संबंधित अनुमान \Elim{\lor} और \Elim{\land} हैं, $!A \ident !B
\lor !C$, और कोई उप-!!{proof}~$\delta_1$ निम्न पर समाप्त होता है:
\[
\AxiomC{}
\RightLabel{$\delta_2$}
\DeduceC{$!D \lor !E$}
\AxiomC{$\Discharge{!D}{x}$}
\RightLabel{$\delta_3$}
\DeduceC{$!B \land !C$}
\AxiomC{$\Discharge{!E}{y}$}
\RightLabel{$\delta_4$}
\DeduceC{$!B \land !C$}
\DischargeRule{\Elim{\lor}}{x\,y}
\TrinaryInfC{$!B \land !C$}
\RightLabel{\Elim{\land}[1]}
\UnaryInfC{$!B$}
\DisplayProof
\]
कट की अंतिम !!{formula}-आवृत्ति~$!A_n$,~\Elim{\lor} का निष्कर्ष, नीचे वाला
$!B \land !C$, है। अंतिम से पहले की !!{formula}-आवृत्ति~$!A_{n-1}$,
~\Elim{\lor} के गौण आधार-वाक्यों में से एक है।

इस कट की लंबाई घटाने के लिए हम \Elim{\lor} और \Elim{\land} अनुमानों का क्रम
बदल (``क्रम-विनिमय कर'') सकते हैं। अर्थात्~$\delta$ में उपप्रमाण~$\delta_1$
को निम्न से बदलते हैं:
\[
\AxiomC{}
\RightLabel{$\delta_2$}
\DeduceC{$!D \lor !E$}
\AxiomC{$\Discharge{!D}{x}$}
\RightLabel{$\delta_3$}
\DeduceC{$!B \land !C$}
\RightLabel{\Elim{\land}[1]}
\UnaryInfC{$!B$}
\AxiomC{$\Discharge{!E}{y}$}
\RightLabel{$\delta_4$}
\DeduceC{$!B \land !C$}
\RightLabel{\Elim{\land}[1]}
\UnaryInfC{$!B$}
\DischargeRule{\Elim{\lor}}{x\,y}
\TrinaryInfC{$!B$}
\DisplayProof
\]
जो कट \Elim{\lor} के नीचे \Elim{\land} के आधार-वाक्य पर समाप्त होते थे, वे
अब \Elim{\lor} के गौण आधार-वाक्यों के ऊपर स्थित \Elim{\land} अनुमानों में से
एक के आधार-वाक्य पर समाप्त होते हैं; अर्थात् प्रत्येक की लंबाई~$1$ घटी है।

इसके अतिरिक्त $\delta_2$, $\delta_3$ या $\delta_4$ के पूर्णतः भीतर, अथवा
~$\delta$ में~$\delta_1$ के पूर्णतः बाहर स्थित कोई भी कट अपरिवर्तित है: उसमें
अब भी वही !!{formula}s और वही अनुमान सम्मिलित हैं, अतः विशेषतः उसकी कोटि और
लंबाई~$\delta$ के संगत कट जैसी हैं। इसलिए नए !!{proof} की कट कोटि नहीं बढ़ी,
और उसकी कट लंबाई पुराने !!{proof} से कम है।

\begin{prob}
क्रम-विनिमय रूपांतरण के बाद कम-से-कम एक कट की लंबाई~$1$ घटती है। यदि हम
\Elim{\land} का~\Elim{\lor} के आर-पार क्रम-विनिमय करते हैं, तो वास्तव में दो
कट खंडों की लंबाई घटती है; इसलिए इस स्थिति में !!{proof} की कट लंबाई
कम-से-कम~$2$ घटती है। क्या ऐसे किसी एक क्रम-विनिमय से !!{proof} की कट लंबाई
~$2$ से अधिक घट सकती है? क्या कट \emph{कोटि} घट सकती है?
\end{prob}

किसी !!a{elimination} नियम को \Elim{\lor} या \Elim{\lexists} नियम के ऊपर ले
जाने की प्रक्रिया को \emph{क्रम-विनिमय रूपांतरण} कहते हैं। कुछ नियम-युग्मों
के प्रतिरूप \olref{tab:permutation-conversions} में दिए गए हैं। (शेष अभ्यास
के लिए छोड़े गए हैं।)

% \begin{table}
% \begin{multline*}
% \shoveleft{\AxiomC{$!D \lor !E$}
% \AxiomC{$\Discharge{!D}{x}$}
% \shortDeduce
% \DeduceC{$!B \land !C$}
% \AxiomC{$\Discharge{!E}{y}$}
% \shortDeduce
% \DeduceC{$!B \land !C$}
% \DischargeRule{\Elim{\lor}}{x\,y}
% \TrinaryInfC{$!B \land !C$}
% \RightLabel{\Elim{\land}[1]}
% \UnaryInfC{$!B$}
% \DisplayProof}\\
% \shoveright{\longrightarrow\ 
% \AxiomC{$!D \lor !E$}
% \AxiomC{$\Discharge{!D}{x}$}
% \shortDeduce
% \DeduceC{$!B \land !C$}
% \RightLabel{\Elim{\land}[1]}
% \UnaryInfC{$!B$}
% \AxiomC{$\Discharge{!E}{y}$}
% \shortDeduce
% \DeduceC{$!B \land !C$}
% \RightLabel{\Elim{\land}[1]}
% \UnaryInfC{$!B$}
% \DischargeRule{\Elim{\lor}}{x\,y}
% \TrinaryInfC{$!B$}
% \DisplayProof}
% \\
% \shoveleft{\AxiomC{$!D \lor !E$}
% \AxiomC{$\Discharge{!D}{x}$}
% \shortDeduce
% \DeduceC{$!B \lif !C$}
% \AxiomC{$\Discharge{!E}{y}$}
% \shortDeduce
% \DeduceC{$!B \lif !C$}
% \DischargeRule{\Elim{\lor}}{x\,y}
% \TrinaryInfC{$!B \lif !C$}
% \AxiomC{$!B$}
% \RightLabel{\Elim{\lif}}
% \BinaryInfC{$!C$}
% \DisplayProof}\\
% \shoveright{\longrightarrow\ 
% \AxiomC{$!D \lor !E$}
% \AxiomC{$\Discharge{!D}{x}$}
% \shortDeduce
% \DeduceC{$!B \lif !C$}
% \AxiomC{$!B$}
% \RightLabel{\Elim{\lif}}
% \BinaryInfC{$!C$}
% \AxiomC{$\Discharge{!E}{y}$}
% \shortDeduce
% \DeduceC{$!B \lif !C$}
% \AxiomC{$!B$}
% \RightLabel{\Elim{\lif}}
% \BinaryInfC{$!C$}
% \DischargeRule{\Elim{\lor}}{x\,y}
% \TrinaryInfC{$!C$}
% \DisplayProof}
% \\
% \shoveleft{\AxiomC{$!D \lor !E$}
% \AxiomC{$\Discharge{!D}{x}$}
% \shortDeduce
% \DeduceC{$!B \lif !C$}
% \AxiomC{$\Discharge{!E}{y}$}
% \shortDeduce
% \DeduceC{$!B \lif !C$}
% \DischargeRule{\Elim{\lor}}{x\,y}
% \TrinaryInfC{$!B \lif !C$}
% \AxiomC{$!B$}
% \RightLabel{\Elim{\lif}}
% \BinaryInfC{$!C$}
% \DisplayProof}\\
% \shoveright{\longrightarrow\ 
% \AxiomC{$!D \lor !E$}
% \AxiomC{$\Discharge{!D}{x}$}
% \shortDeduce
% \DeduceC{$!B \lif !C$}
% \AxiomC{$!B$}
% \RightLabel{\Elim{\lif}}
% \BinaryInfC{$!C$}
% \AxiomC{$\Discharge{!E}{y}$}
% \shortDeduce
% \DeduceC{$!B \lif !C$}
% \AxiomC{$!B$}
% \RightLabel{\Elim{\lif}}
% \BinaryInfC{$!C$}
% \DischargeRule{\Elim{\lor}}{x\,y}
% \TrinaryInfC{$!C$}
% \DisplayProof}
% \\
% \shoveleft{\AxiomC{$!D \lor !E$}
% \AxiomC{$\Discharge{!D}{x}$}
% \shortDeduce
% \DeduceC{$\lforall[x][!B(x)]$}
% \AxiomC{$\Discharge{!E}{y}$}
% \shortDeduce
% \DeduceC{$\lforall[x][!B(x)]$}
% \DischargeRule{\Elim{\lor}}{x\,y}
% \TrinaryInfC{$\lforall[x][!B(x)]$}
% \RightLabel{\Elim{\lforall}}
% \UnaryInfC{$!B(t)$}
% \DisplayProof}\\
% \shoveright{\longrightarrow\ 
% \AxiomC{$!D \lor !E$}
% \AxiomC{$\Discharge{!D}{x}$}
% \shortDeduce
% \DeduceC{$\lforall[x][!B(x)]$}
% \RightLabel{\Elim{\lforall}}
% \UnaryInfC{$!B(t)$}
% \AxiomC{$\Discharge{!E}{y}$}
% \shortDeduce
% \DeduceC{$\lforall[x][!B(x)]$}
% \RightLabel{\Elim{\lforall}}
% \UnaryInfC{$!B(t)$}
% \DischargeRule{\Elim{\lor}}{x\,y}
% \TrinaryInfC{$!B(t)$}
% \DisplayProof}
% \end{multline*}
% \caption{Some permutation conversions for \Log{N1i}.}
% \ollabel{tab:permutation-conversions}
% \end{table}

{\small
\begin{longtable}{@{}l@{}}
\AxiomC{$!D \lor !E$}
\AxiomC{$\Discharge{!D}{x}$}
\shortDeduce
\DeduceC{$!B \land !C$}
\AxiomC{$\Discharge{!E}{y}$}
\shortDeduce
\DeduceC{$!B \land !C$}
\DischargeRule{\Elim{\lor}}{x\,y}
\TrinaryInfC{$!B \land !C$}
\RightLabel{\Elim{\land}[1]}
\UnaryInfC{$!B$}
\DisplayProof
$\longrightarrow$\\[6ex]
\quad\AxiomC{$!D \lor !E$}
\AxiomC{$\Discharge{!D}{x}$}
\shortDeduce
\DeduceC{$!B \land !C$}
\RightLabel{\Elim{\land}[1]}
\UnaryInfC{$!B$}
\AxiomC{$\Discharge{!E}{y}$}
\shortDeduce
\DeduceC{$!B \land !C$}
\RightLabel{\Elim{\land}[1]}
\UnaryInfC{$!B$}
\DischargeRule{\Elim{\lor}}{x\,y}
\TrinaryInfC{$!B$}
\DisplayProof
\\[10ex]
\AxiomC{$!D \lor !E$}
\AxiomC{$\Discharge{!D}{x}$}
\shortDeduce
\DeduceC{$!B \lif !C$}
\AxiomC{$\Discharge{!E}{y}$}
\shortDeduce
\DeduceC{$!B \lif !C$}
\DischargeRule{\Elim{\lor}}{x\,y}
\TrinaryInfC{$!B \lif !C$}
\AxiomC{$!B$}
\RightLabel{\Elim{\lif}}
\BinaryInfC{$!C$}
\DisplayProof
$\longrightarrow$\\[6ex]
\AxiomC{$!D \lor !E$}
\AxiomC{$\Discharge{!D}{x}$}
\shortDeduce
\DeduceC{$!B \lif !C$}
\AxiomC{$!B$}
\RightLabel{\Elim{\lif}}
\BinaryInfC{$!C$}
\AxiomC{$\Discharge{!E}{y}$}
\shortDeduce
\DeduceC{$!B \lif !C$}
\AxiomC{$!B$}
\RightLabel{\Elim{\lif}}
\BinaryInfC{$!C$}
\DischargeRule{\Elim{\lor}}{x\,y}
\TrinaryInfC{$!C$}
\DisplayProof
\\[10ex]
\AxiomC{$!D \lor !E$}
\AxiomC{$\Discharge{!D}{x}$}
\shortDeduce
\DeduceC{$!B \lif !C$}
\AxiomC{$\Discharge{!E}{y}$}
\shortDeduce
\DeduceC{$!B \lif !C$}
\DischargeRule{\Elim{\lor}}{x\,y}
\TrinaryInfC{$!B \lif !C$}
\AxiomC{$!B$}
\RightLabel{\Elim{\lif}}
\BinaryInfC{$!C$}
\DisplayProof
$\longrightarrow$\\[6ex]
\AxiomC{$!D \lor !E$}
\AxiomC{$\Discharge{!D}{x}$}
\shortDeduce
\DeduceC{$!B \lif !C$}
\AxiomC{$!B$}
\RightLabel{\Elim{\lif}}
\BinaryInfC{$!C$}
\AxiomC{$\Discharge{!E}{y}$}
\shortDeduce
\DeduceC{$!B \lif !C$}
\AxiomC{$!B$}
\RightLabel{\Elim{\lif}}
\BinaryInfC{$!C$}
\DischargeRule{\Elim{\lor}}{x\,y}
\TrinaryInfC{$!C$}
\DisplayProof
\\[10ex]
\AxiomC{$!D \lor !E$}
\AxiomC{$\Discharge{!D}{x}$}
\shortDeduce
\DeduceC{$\lforall[x][!B(x)]$}
\AxiomC{$\Discharge{!E}{y}$}
\shortDeduce
\DeduceC{$\lforall[x][!B(x)]$}
\DischargeRule{\Elim{\lor}}{x\,y}
\TrinaryInfC{$\lforall[x][!B(x)]$}
\RightLabel{\Elim{\lforall}}
\UnaryInfC{$!B(t)$}
\DisplayProof
\quad$\longrightarrow$\\[8ex]
\quad\AxiomC{$!D \lor !E$}
\AxiomC{$\Discharge{!D}{x}$}
\shortDeduce
\DeduceC{$\lforall[x][!B(x)]$}
\RightLabel{\Elim{\lforall}}
\UnaryInfC{$!B(t)$}
\AxiomC{$\Discharge{!E}{y}$}
\shortDeduce
\DeduceC{$\lforall[x][!B(x)]$}
\RightLabel{\Elim{\lforall}}
\UnaryInfC{$!B(t)$}
\DischargeRule{\Elim{\lor}}{x\,y}
\TrinaryInfC{$!B(t)$}
\DisplayProof
\\
\caption{\Log{N1i} के कुछ क्रम-विनिमय रूपांतरण।}
\ollabel{tab:permutation-conversions}
\end{longtable}
}

\begin{prob}
\Elim{\lor} के बाद आने वाले \Elim{\lor} या \Elim{\lnot} के लिए क्रम-विनिमय
रूपांतरण दीजिए।
\end{prob}

\begin{prob}
\Elim{\lexists} के बाद आने वाले \Elim{\exists} के लिए क्रम-विनिमय रूपांतरण
दीजिए। समझाइए कि बाद वाली स्थिति में आइगेनचर की किसी शर्त का उल्लंघन क्यों
नहीं होता।
\end{prob}

हमें आश्वस्त कर लेना चाहिए कि क्रम-विनिमय रूपांतरण लगाते समय हम कहीं आइगेनचर
की किसी शर्त का उल्लंघन तो नहीं करते। विचार करें
\[
\AxiomC{}
\RightLabel{$\delta_2$}
\DeduceC{$!D \lor !E$}
\AxiomC{$\Discharge{!D}{x}$}
\RightLabel{$\delta_3$}
\DeduceC{$\lexists[x][!B(x)]$}
\AxiomC{$\Discharge{!E}{y}$}
\RightLabel{$\delta_4$}
\DeduceC{$\lexists[x][!B(x)]$}
\DischargeRule{\Elim{\lor}}{x\,y}
\TrinaryInfC{$\lexists[x][!B(x)]$}
\AxiomC{$\Discharge{!B(c)}{z}$}
\RightLabel{$\delta_5$}
\DeduceC{$!C$}
\DischargeRule{\Elim{\lexists}}{z}
\BinaryInfC{$!C$}
\DisplayProof
\]
यदि हम इसे निम्न से बदलें
\[
\AxiomC{}
\RightLabel{$\delta_2$}
\DeduceC{$!D \lor !E$}
\AxiomC{$\Discharge{!D}{x}$}
\RightLabel{$\delta_3$}
\DeduceC{$\lexists[x][!B(x)]$}
\AxiomC{$\Discharge{!B(c)}{z}$}
\RightLabel{$\delta_5$}
\DeduceC{$!C$}
\DischargeRule{\Elim{\lexists}}{z}
\BinaryInfC{$!C$}
\AxiomC{$\Discharge{!E}{y}$}
\RightLabel{$\delta_4$}
\DeduceC{$\lexists[x][!B(x)]$}
\AxiomC{$\Discharge{!B(c)}{z}$}
\RightLabel{$\delta_5$}
\DeduceC{$!C$}
\DischargeRule{\Elim{\lexists}}{z}
\BinaryInfC{$!C$}
\DischargeRule{\Elim{\lor}}{x\,y}
\TrinaryInfC{$!C$}
\DisplayProof
\]
तो हमें चिंता हो सकती है कि आइगेनचर~$c$, मान लें,~$!D$ में आता हो। आखिर
मान्यता $!D$ मूल~$\Elim{\lexists}$ के ऊपर मुक्त होती है, इसलिए मूल !!{proof}
में वह ऐसी मान्यता में नहीं आती जो~\Elim{\exists} अनुमान के मुख्य आधार-वाक्य
के ऊपर खुली हो; अतः वहाँ आइगेनचर की शर्त पूरी होती है। पर नए !!{proof} में
~$!D$,~\Elim{\lexists} अनुमानों के बाद तक मुक्त नहीं होती, इसलिए आइगेनचर की
शर्त का उल्लंघन होता प्रतीत होगा। किंतु याद रखें कि हम अपने !!{proof}s को
नियमित मानते हैं: प्रत्येक आइगेनचर केवल एक अनुमान का आइगेनचर है और !!{proof}
में केवल \Elim{\lexists} के गौण आधार-वाक्य के ऊपर आता है। विशेषतः $c$,
$\delta_3$ या~$\delta_4$ में नहीं आ सकता।

यह उदाहरण एक और बात दिखाता है: ध्यान दें कि नए !!{proof} में~$\delta_5$ की
दो प्रतियाँ हैं। इसका अर्थ है कि यदि $\delta_5$ में कोई महत्तम कट है, तो हमने
कट लंबाई \emph{नहीं} घटाई! इसलिए हमें क्रम-विनिमय रूपांतरण केवल तभी लगाना
चाहिए जब हमें ज्ञात हो कि $\delta_5$ में कोई महत्तम कट नहीं है। इसे हम सदा
किसी \emph{सबसे ऊपर के} महत्तम कट को चुनकर सुनिश्चित करेंगे; अर्थात् ऐसा
महत्तम कट जिसके उप-!!{proof} (यहाँ~$\delta_1$) में उप-!!{proof} के अंतिम अनुमान
के ऊपर समाप्त होने वाला कोई महत्तम कट न हो। इससे~$\delta_5$ में (और वास्तव
में $\delta_2$, $\delta_3$ या~$\delta_4$ में भी) अन्य महत्तम कट होने की
संभावना समाप्त हो जाती है।

\end{document}
